\begin{abstract}
Language models can obey a literal instruction not to reveal a secret while still letting that secret shape an unrelated story. We test whether an external outline reduces this nonliteral leakage by offloading future content decisions, and whether secret-aligned residual activity explains any reduction. Using a local \gemma writer and judge, we cross 15 published secret words with matched no-secret controls, two neutral story briefs, and four context conditions. We also run an unconstrained anchor, a six-item plot-twist pilot, residual-stream measurements, and projection-removal interventions. The unconstrained anchor replicates leakage: secret-aware stories receive a 16.33-point higher order-balanced influence score than controls (95\% bootstrap CI $[4.73,29.10]$, $p=.020$). A neutral brief removes the detectable signal before an outline is added (no-outline mean $-0.15$, CI $[-2.87,2.40]$). Consequently, the outline contrast is inconclusive: outline minus no outline is $-2.50$ (CI $[-6.83,1.85]$, $p=.280$), and outline minus length-matched irrelevant context is $-2.08$ (CI $[-7.12,2.95]$, $p=.418$). Carrier-defined residual gaps are small and inconsistent. Projection removal strongly inverts judgments but shortens stories from 488 to 102 words on average, invalidating a selective causal interpretation. These results identify freedom over story content as a critical moderator of detectable leakage, but do not establish planning offload or a selectively causal linear secret direction.
\end{abstract}

